% math0 manuscript. Build with: latexmk -pdf paper.tex
\documentclass[11pt]{amsart}
\usepackage[margin=1.1in]{geometry}
\usepackage{amsmath,amssymb,amsthm,mathtools}
\usepackage{url}
\usepackage{hyperref}
\usepackage{xcolor}
\usepackage{listings}
\usepackage{newunicodechar}

\lstdefinelanguage{lean}{
  keywords={theorem,lemma,def,abbrev,example,fun,by,have,show,from,let,in,if,then,else,match,with,
    namespace,end,open,import,section,variable,instance,structure,where,calc,obtain,intro,exact,
    refine,rcases,rintro,simp,omega,decide,norm_num,linarith,nlinarith,positivity,induction,cases,
    private},
  sensitive=true, comment=[l]{--}, morecomment=[s]{/-}{-/}, string=[b]"}
\lstset{language=lean, basicstyle=\ttfamily\footnotesize, columns=fullflexible, keepspaces=true,
  breaklines=true, keywordstyle=\color{blue!60!black}\bfseries, commentstyle=\color{gray},
  frame=single, extendedchars=true, inputencoding=utf8, upquote=true,
  literate={ℕ}{{$\mathbb{N}$}}1 {∀}{{$\forall$}}1 {∃}{{$\exists$}}1 {∧}{{$\wedge$}}1
    {→}{{$\to$}}1 {≤}{{$\le$}}1 {∣}{{$\mid$}}1 {×}{{$\times$}}1}
% Unicode commonly found in Lean sources.
\newunicodechar{ℕ}{\ensuremath{\mathbb{N}}}
\newunicodechar{ℤ}{\ensuremath{\mathbb{Z}}}
\newunicodechar{ℚ}{\ensuremath{\mathbb{Q}}}
\newunicodechar{ℝ}{\ensuremath{\mathbb{R}}}
\newunicodechar{ℂ}{\ensuremath{\mathbb{C}}}
\newunicodechar{∀}{\ensuremath{\forall}}
\newunicodechar{∃}{\ensuremath{\exists}}
\newunicodechar{→}{\ensuremath{\to}}
\newunicodechar{←}{\ensuremath{\leftarrow}}
\newunicodechar{↔}{\ensuremath{\leftrightarrow}}
\newunicodechar{↦}{\ensuremath{\mapsto}}
\newunicodechar{¬}{\ensuremath{\neg}}
\newunicodechar{∧}{\ensuremath{\wedge}}
\newunicodechar{∨}{\ensuremath{\vee}}
\newunicodechar{≤}{\ensuremath{\le}}
\newunicodechar{≥}{\ensuremath{\ge}}
\newunicodechar{≠}{\ensuremath{\ne}}
\newunicodechar{∈}{\ensuremath{\in}}
\newunicodechar{∉}{\ensuremath{\notin}}
\newunicodechar{⊆}{\ensuremath{\subseteq}}
\newunicodechar{∩}{\ensuremath{\cap}}
\newunicodechar{∪}{\ensuremath{\cup}}
\newunicodechar{∑}{\ensuremath{\sum}}
\newunicodechar{∏}{\ensuremath{\prod}}
\newunicodechar{⟨}{\ensuremath{\langle}}
\newunicodechar{⟩}{\ensuremath{\rangle}}
\newunicodechar{·}{\ensuremath{\cdot}}
\newunicodechar{∘}{\ensuremath{\circ}}
\newunicodechar{⁻}{\ensuremath{^{-}}}
\newunicodechar{¹}{\ensuremath{^{1}}}
\newunicodechar{²}{\ensuremath{^{2}}}
\newunicodechar{₀}{\ensuremath{_{0}}}
\newunicodechar{₁}{\ensuremath{_{1}}}
\newunicodechar{₂}{\ensuremath{_{2}}}
\newunicodechar{α}{\ensuremath{\alpha}}
\newunicodechar{β}{\ensuremath{\beta}}
\newunicodechar{ε}{\ensuremath{\varepsilon}}
\newunicodechar{δ}{\ensuremath{\delta}}
\newunicodechar{π}{\ensuremath{\pi}}
\newunicodechar{σ}{\ensuremath{\sigma}}
\newunicodechar{φ}{\ensuremath{\varphi}}
\newunicodechar{∞}{\ensuremath{\infty}}
\newunicodechar{⌊}{\ensuremath{\lfloor}}
\newunicodechar{⌋}{\ensuremath{\rfloor}}
\newunicodechar{⌈}{\ensuremath{\lceil}}
\newunicodechar{⌉}{\ensuremath{\rceil}}
\newunicodechar{∣}{\ensuremath{\mid}}
\newunicodechar{≡}{\ensuremath{\equiv}}
\newunicodechar{×}{\ensuremath{\times}}
\newunicodechar{⧸}{\ensuremath{/}}
\newunicodechar{‖}{\ensuremath{\|}}
\newunicodechar{√}{\ensuremath{\surd}}
\newunicodechar{∅}{\ensuremath{\emptyset}}
\newunicodechar{⦃}{\ensuremath{\{\!|}}
\newunicodechar{⦄}{\ensuremath{|\!\}}}
\newunicodechar{𝓝}{\ensuremath{\mathcal{N}}}
\newunicodechar{≫}{\ensuremath{\gg}}
\newunicodechar{≪}{\ensuremath{\ll}}

\theoremstyle{plain}
\newtheorem{theorem}{Theorem}
\newtheorem{lemma}[theorem]{Lemma}
\newtheorem{proposition}[theorem]{Proposition}
\newtheorem{corollary}[theorem]{Corollary}
\theoremstyle{definition}
\newtheorem{definition}[theorem]{Definition}
\newtheorem{remark}[theorem]{Remark}

\emergencystretch=3em
% Lean identifiers in running text (breakable at dots and underscores).
\expandafter\def\expandafter\UrlBreaks\expandafter{\UrlBreaks\do\_\do\.}
\DeclareUrlCommand\decl{\urlstyle{tt}}
\newcommand{\N}{\mathbb{N}}
\newcommand{\rf}[2]{#1^{\overline{#2}}}
\newcommand{\fl}[1]{\lfloor #1\rfloor}

\title[The divisible-product triangle A090905 and the Bertrand primes]{The rows of the divisible-product triangle OEIS A090905 end at twice the Bertrand primes}
\author{math0}
\address{Produced by Claude (an AI model developed by Anthropic) in the math0 project}
\date{October 20, 2026}

\begin{document}
\begin{abstract}
Split $1,2,3,\dots$ into consecutive blocks: block $1$ is $(1)$, and each later block is the shortest block, not shorter than its predecessor, whose product is a multiple of the product of its predecessor (OEIS A090905).
We prove Amarnath Murthy's 2003 conjecture from this entry: for every $n>4$ the last term of block $n$ is $2p$, where $p$ is the largest prime of block $n-1$, and $p$ is the Bertrand prime $B_{n-1}$ (OEIS A006992: $B_1=2$, and $B_{k+1}$ is the largest prime below $2B_k$).
The proof combines Legendre's formula with a prime in $(5x/3,2x)$ for every integer $x\ge83$, a weak form of Nagura's 1952 theorem, which we reprove by an explicit prime ladder and Chebyshev's method.
We know of no earlier proof, so we classify the result as new, although the argument is elementary once the prime-gap bound is available.
The theorem is verified in Lean~4 against the frozen, fidelity-reviewed math0 formalization of the conjecture, whose definitions are checked by the kernel against the OEIS data.
\end{abstract}
\maketitle

\section{Introduction}

OEIS A090905 \cite{A090905}, contributed by Amarnath Murthy on December~13, 2003, is the ``Left side of irregular triangle of natural numbers in which every row product is a multiple of the previous''. The triangle splits $1,2,3,\dots$ into consecutive blocks (``groups'') $(1)$, $(2)$, $(3,4)$, $(5,\dots,8)$, $(9,\dots,14)$, $(15,\dots,26)$, $(27,\dots,46)$, \dots\ Each block starts right after the previous one, and its length is the least $k$, not smaller than the previous length, such that the product of the previous block divides the product of the new one; this is the rule of the Mathematica program of M.~De Vlieger (2016) in the entry. The companion entries A090904, A090906 and A090907 \cite{A090904} list the row products, the row lengths, and the ratios of consecutive row products. All four entries carry Murthy's conjecture (quoted in Section~\ref{sec:statement}): for $n>4$ the last term of row $n$ is twice the largest prime of row $n-1$, and these primes are the Bertrand primes A006992 \cite{A006992}. We prove it (Theorem~\ref{thm:main}).

\subsection*{Method}
Let $b<a<c$ be consecutive Bertrand primes and suppose row $n$ is $[2b+1,2a]$. The prime $c$ lies in row $n$, and no multiple of $c$ lies in $(2a,2c)$, so row $n+1$ must reach $2c$; ending at $2c$ suffices because $(2a)!^2\mid(2b)!\,(2c)!$ (Legendre's formula). Both steps use $3a>5b$ and $3c>5a$, i.e.\ a prime in $(5x/3,2x)$ for every integer $x\ge83$. This follows from Nagura's theorem \cite{Nag} (a prime in $(x,6x/5)$ for $x\ge25$); the Lean development proves the needed form from scratch, by a ladder of primes below $10^6$ and by Chebyshev's method \cite{Cheb} above.

\subsection*{Status and prior work}
The entries A090904--A090907 (last edited in December 2016) state the claim only as a conjecture. A090907 \cite{A090904} gives a formula for the ratios in terms of A006992, ``verified for $4<n<21$''. A090906 \cite{A090904} contains an unproved formula of M.~Bouayoun (March 2004) for the row lengths, $2(\mathrm{A006992}(n)-\mathrm{A006992}(n-1))$ for $n>4$; as written it is off by one index (it gives $12$ for $n=5$, whereas the fifth row length is $6$). Theorem~\ref{thm:main} gives the shifted form $2(\mathrm{A006992}(n-1)-\mathrm{A006992}(n-2))$ for $n\ge6$, and the formula of A090907 for $n\ge6$ (neither consequence is formalized). A381901 \cite{A381901} (A.~Niedermaier, March 2025) defines the same partition without the lower bound on the row length and states, without proof, that the last term of row $n$ is $2\,\mathrm{A006992}(n-1)$ for $n\ge5$; this follows from Theorem~\ref{thm:main} and Remark~\ref{rem:bound}. Related entries are A113117, A113118 (L.~Quet, 2006) and Murthy's variant A093911 \cite{A093911}. The math0 referees found no proof or counterexample in searches of arXiv, MathOverflow, Mathematics Stack Exchange and SeqFan-style queries, and no formalization in the formal-conjectures repository \cite{FC}.

\section{Statement}\label{sec:statement}

For integers $s\ge1$ and $k\ge0$ write $\rf{s}{k}=s(s+1)\cdots(s+k-1)$ (Lean's \decl{Nat.ascFactorial}), the product of the block $s,\dots,s+k-1$.

\begin{definition}\label{def:tri}
Let $a(1)=1$, $L(1)=1$ and, for $n\ge1$,
\[
a(n+1)=a(n)+L(n),\qquad L(n+1)=\min\{k\ge L(n): \rf{a(n)}{L(n)}\mid \rf{a(n+1)}{k}\}.
\]
Row $n$ is the block $[a(n),e(n)]$ with $e(n)=a(n)+L(n)-1$. The Bertrand primes are $B_1=2$ and $B_{k+1}=$ the largest prime $<2B_k$; thus $(B_k)_{k\ge1}=2,3,5,7,13,23,43,83,163,317,\dots$
\end{definition}

The minimum exists: $k=a(n+1)-1=a(n)+L(n)-1\ge L(n)$ is admissible, because the row product divides $(a(n+1)-1)!$, which divides the product of any $a(n+1)-1$ consecutive integers. The COMMENTS section of A090905 reads, verbatim (the entry is by Amarnath Murthy, December~13, 2003):
\begin{quote}
{\small\raggedright\texttt{Conjecture: For n > 4 the last term of the n-th group is 2p where p is the largest prime in the (n-1)th group. And these are the Bertrand primes.}\par}
\end{quote}

\begin{theorem}\label{thm:main}
For every $n\ge5$, the set of primes in row $n-1$ has a largest element $p$; moreover $e(n)=2p$ and $p=B_{n-1}$.
\end{theorem}

Applying this to $n$ and to $n-1$ and using $a(n)=e(n-1)+1$, we see that row $n$ is $[2B_{n-2}+1,\,2B_{n-1}]$ and $L(n)=2B_{n-1}-2B_{n-2}$ for every $n\ge6$.

\subsection*{Lean formalization}
The statement was formalized by math0 in \path{lean/Math0/Statements/OEIS/A090905.lean}, reviewed for fidelity and frozen (SHA-256 \texttt{8616fe97\ldots a42a44}) before any proof attempt. Its definitions read as follows (verbatim; docstrings and the auxiliary theorems are omitted):
\begin{lstlisting}
namespace Math0.Statements.A090905
def search (P s : ℕ) : ℕ → ℕ → ℕ → ℕ
  | 0, k, _ => k
  | fuel + 1, k, Q => if P ∣ Q then k else search P s fuel (k + 1) ((s + k) * Q)
def next (r : ℕ × ℕ) : ℕ × ℕ :=
  (r.1 + r.2,
    search (r.1.ascFactorial r.2) (r.1 + r.2) (r.1 + r.2) r.2 ((r.1 + r.2).ascFactorial r.2))
def row (n : ℕ) : ℕ × ℕ := next^[n - 1] (1, 1)
def a (n : ℕ) : ℕ := (row n).1
def len (n : ℕ) : ℕ := (row n).2
def last (n : ℕ) : ℕ := a n + len n - 1
def bertrand (n : ℕ) : ℕ := (fun b => Nat.findGreatest Nat.Prime (2 * b - 1))^[n - 1] 2
def statement : Prop :=
  ∀ n : ℕ, 4 < n → ∃ p : ℕ,
    IsGreatest {q : ℕ | q.Prime ∧ a (n - 1) ≤ q ∧ q ≤ last (n - 1)} p ∧
      last n = 2 * p ∧ p = bertrand (n - 1)
\end{lstlisting}
A row is stored as the pair (first term, length). Called with $Q=\rf{s}{k}$, \lstinline|search P s fuel k Q| returns the least $j\in[k,k+\mathit{fuel})$ with $P\mid\rf{s}{j}$ (or $k+\mathit{fuel}$ if there is none). The file proves \decl{a_one}, \decl{len_one}, \decl{a_succ} and \decl{len_succ_isLeast}, i.e.\ the recursion of Definition~\ref{def:tri} for $n\ge1$; in particular the fuel $a(n+1)$ always suffices (the admissible $k=a(n+1)-1$ lies in the search window), so \decl{a}, \decl{len}, \decl{last} are exactly $a$, $L$, $e$ (and $a(n)\ge1$, so the truncated subtraction in \decl{last} is harmless). Since $B_k\ge2$, \lstinline|Nat.findGreatest Nat.Prime (2 * b - 1)| at $b=B_k$ is the largest prime $<2B_k$, so \decl{bertrand} is A006992. Lemmas proved by \texttt{decide +kernel} check $a(1),\dots,a(18)$, $L(1),\dots,L(17)$ and $B_1,\dots,B_{13}$ against the data of A090905, A090906 and A006992 (\decl{a_data}, \decl{len_data}, \decl{bertrand_data}).

``The largest prime in the $(n-1)$th group'' is the greatest element of the set of primes in $[a(n-1),e(n-1)]$, whose existence is asserted. ``These are the Bertrand primes'' is read as $p=B_{n-1}$: for $n=5$ the largest prime of row $4=(5,\dots,8)$ is $7=B_4$, and the rows are disjoint and increasing ($2,3,5$ do not occur as such $p$). This is the strongest natural reading; the weaker reading ``each such $p$ is a Bertrand prime'' follows from it.

\begin{remark}\label{rem:bound}
The bound $k\ge L(n)$ comes from De Vlieger's program, not from the entry's name. Dropping it does not change the triangle (not formalized): rows $1$--$20$ agree by computer check, and for $n\ge10$ the minimality argument of Lemma~\ref{lem:step} does not use the bound.
\end{remark}

\section{Proof}\label{sec:proof}

We follow the Lean proof; names of its (private) lemmas are given in brackets.

\subsection{\texorpdfstring{A prime between $5x/3$ and $2x$}{A prime between 5x/3 and 2x}}

\begin{proposition}[\decl{exists_prime}]\label{prop:nagura}
For every integer $x\ge83$ there is a prime $p$ with $5x<3p<6x$.
\end{proposition}

This follows from Nagura's theorem \cite{Nag} at $5x/3\ge25$; we give the self-contained proof that is formalized.

\begin{lemma}[\decl{ladder_aux}]\label{lem:ladder}
Let $q_0=138$, and let $q_1<\dots<q_{53}$ be the primes $163$, $193$, $229$, \dots, $1491013$, $1789201$ of the list \decl{ladder}, which satisfy $5q_{i}<6q_{i-1}$ for $1\le i\le53$. Then Proposition~\ref{prop:nagura} holds for $83\le x<10^6$.
\end{lemma}

\begin{proof}
We have $3q_0=414<5x$ and $5x<5\cdot10^6<3q_{53}=5367603$. Let $i\ge1$ be least with $5x<3q_i$. Then $3q_{i-1}\le5x$, so $15q_i<18q_{i-1}\le30x$, i.e.\ $3q_i<6x$. In Lean, \texttt{norm\_num} produces a primality certificate for each $q_i$, which the kernel checks, and the $53$ inequalities $5q_i<6q_{i-1}$ are checked by \texttt{decide}.
\end{proof}

Let $\Lambda$ be the von Mangoldt function, $\psi(y)=\sum_{d\le y}\Lambda(d)$ and $\theta(y)=\sum_{p\le y}\log p$.

The identity $\log M!=\sum_{d\le M}\Lambda(d)\fl{M/d}$ follows by induction on $M$ from $\fl{(M+1)/d}-\fl{M/d}=[d\mid M+1]$ and $\sum_{d\mid n}\Lambda(d)=\log n$; with $\fl{\fl{N/k}/d}=\fl{\fl{N/d}/k}$ it gives $\log\fl{N/k}!=\sum_{d\le N}\Lambda(d)\fl{\fl{N/d}/k}$ for $k\ge1$ (\decl{log_factorial_div_eq_sum}). A \emph{weight} is a list $W$ of pairs $(c_i,k_i)\in\mathbb{Z}\times\N_{\ge1}$. Put $H_W(n)=\sum_i c_i\fl{n/k_i}$ and $G_W(N)=\sum_i c_i\log\fl{N/k_i}!$. Then
\begin{equation}\label{eq:G}
G_W(N)=\sum_{d\le N}\Lambda(d)\,H_W(\fl{N/d}).
\end{equation}

\begin{lemma}[\decl{Gw_le_psi}, \decl{psi_le_Gw}]\label{lem:weights}
\textup{(a)} If $H_W\le1$ on $\N$, then $G_W(N)\le\psi(N)$.
\textup{(b)} If $H_W\ge0$ on $\N$ and $H_W(n)\ge1$ whenever $P\nmid n$, then $\psi(N)\le G_W(N)+\psi(\fl{N/P})$.
\end{lemma}

\begin{proof}
(a) is immediate from \eqref{eq:G} and $\Lambda\ge0$. For (b), split $\psi(N)=\sum_{d\le N}\Lambda(d)$ according to whether $P\mid\fl{N/d}$. If so, then $P\le\fl{N/d}$ (as $\fl{N/d}\ge1$), hence $d\le\fl{N/P}$, and these terms contribute at most $\psi(\fl{N/P})$. The other terms satisfy $\Lambda(d)\le\Lambda(d)H_W(\fl{N/d})$, and the remaining terms of \eqref{eq:G} are $\ge0$.
\end{proof}

We use Chebyshev's weight \cite{Cheb} and a $19$-term weight found by linear programming over the divisors of $420$:
\begin{align*}
W_L&=(1,1),(-1,2),(-1,3),(-1,5),(1,30),\\
W_U&=(1,1),(-1,2),(-1,3),(-1,5),(1,6),(-1,7),(1,10),(-1,12),(1,14),(-1,20),\\
&\quad\ (-1,30),(1,35),(-1,42),(1,60),(-1,70),(1,84),(-1,105),(1,210),(-4,420).
\end{align*}
If every $k_i$ divides $P$, then $H_W(n)=H_W(n\bmod P)+\fl{n/P}\,R_W(P)$ with $R_W(P)=\sum_ic_iP/k_i$ (\decl{Hw_mod}). Here $R_{W_L}(30)=R_{W_U}(420)=0$, so $H_{W_L}$ has period $30$ and $H_{W_U}$ has period $420$. Evaluating one period in the kernel (\texttt{decide} in \decl{HL_le}, \texttt{decide +kernel} in \decl{HU_spec}) gives $H_{W_L}\le1$, $H_{W_U}\ge0$, and $H_{W_U}(n)\ge1$ for $420\nmid n$. Lemma~\ref{lem:weights} then gives $G_{W_L}(N)\le\psi(N)\le G_{W_U}(N)+\psi(\fl{N/420})$.

\begin{lemma}[\decl{log_fact_lower}, \decl{log_fact_upper}]\label{lem:stirling}
For $M\ge1$: $M\log M-M+\tfrac12\log M+\tfrac12\le\log M!\le M\log M-M+\tfrac12\log M+1$.
\end{lemma}

\begin{proof}
The lower bound is \decl{Stirling.le_log_factorial_stirling} with $\log(2\pi)\ge1$. For the upper bound, $M!/(\sqrt{2M}\,(M/e)^M)$ is non-increasing in $M\ge1$ (\decl{Stirling.log_stirlingSeq'_antitone}) and equals $e/\sqrt2$ at $M=1$.
\end{proof}

\begin{lemma}[\decl{Gw_bounds}, \decl{psi_lower}, \decl{psi_upper}]\label{lem:psi}
For every integer $m\ge1$, \ $386.94\,m-5.025-\tfrac12\log m\le\psi(420m)\le448.99\,m-3.685-3\log m+\psi(m)$.
\end{lemma}

\begin{proof}
Let $W$ be $W_L$ or $W_U$ and $r_i=420/k_i\in\N$, so that $\fl{420m/k_i}=r_im$. Apply Lemma~\ref{lem:stirling} to $M=r_im$, write $\log M=\log r_i+\log m$, multiply by $c_i$ (reversing the inequalities when $c_i<0$) and sum. Since $\sum_ic_ir_i=0$, the terms in $m\log m$ and $m$ cancel, and
\[
mD_1+\tfrac12D_2+\tfrac S2\log m+E^-\ \le\ G_W(420m)\ \le\ mD_1+\tfrac12D_2+\tfrac S2\log m+E^+,
\]
where $D_1=\sum c_ir_i\log r_i$, $D_2=\sum c_i\log r_i$, $S=\sum c_i$, $E^+=\sum_{c_i>0}c_i+\sum_{c_i<0}c_i/2$ and $E^-=\sum_{c_i>0}c_i/2+\sum_{c_i<0}c_i$. For $W_L$: $D_1=196\log2+126\log3+70\log5\ge386.94$, $D_2=-\log420\ge-6.05$, $S=-1$, $E^-=-2$. For $W_U$: $D_1=192\log2+123\log3+70\log5+35\log7\le448.99$, $D_2=-\log11760\le-9.37$, $S=-6$, $E^+=1$. These numerical bounds use Mathlib's $0.6931471803<\log2<0.6931471808$ and $16$ terms of the series of $\log(1-t)$ at $t=\frac14,\frac15,\frac18$ with Mathlib's remainder bound \decl{Real.abs_log_sub_add_sum_range_le}. Combining with $G_{W_L}(420m)\le\psi(420m)\le G_{W_U}(420m)+\psi(m)$ gives the claim.
\end{proof}

\begin{proof}[Proof of Proposition~\ref{prop:nagura} for $x\ge10^6$ \textup{(\decl{exists_prime_large})}]
For an integer $p$, $5x<3p<6x$ means $\fl{5x/3}<p\le2x-1$. Suppose there is no such prime; then $\theta(2x-1)=\theta(\fl{5x/3})$. Let $m_0=\fl{(2x-1)/420}$, $m_1=\fl{\fl{5x/3}/420}+1$, $s=\sqrt x\ge1000$ and $\lambda=\log x$. Then $1\le m_0,m_1\le x$, $420m_0\le2x-1$, $2x\le420m_0+420$, $\fl{5x/3}\le420m_1$, $1260m_1\le5x+1260$ and $196m_1\le x$. We use Mathlib's bounds $\psi(y)-\theta(y)\le2\sqrt y\log y$ and $\psi(y)\le y\log4+2\sqrt y\log y$ for $y\ge1$, the monotonicity of $\psi$, $\theta\le\psi$, Lemma~\ref{lem:psi}, $\sqrt{2x-1}\le1.4143\,s$, $\log(2x-1)\le\log2+\lambda$, $0\le\log m_i\le\lambda$ and $\sqrt{m_1}\le s/14$:
\begin{align*}
\theta(2x-1)&\ge\psi(420m_0)-2\sqrt{2x-1}\log(2x-1)\ge\frac{386.94\,x}{210}-391.965-\frac\lambda2-2.8286\,s(\log2+\lambda),\\
\theta(\fl{5x/3})&\le\psi(420m_1)\le\frac{448.99\,x}{252}+445.305+\frac{1.3864\,x}{196}+\frac{s\lambda}{7}.
\end{align*}
The left sides are equal, $386.94/210-448.99/252-1.3864/196>0.05379$ and $2.8286\log2\le1.9608$, so $0.05379\,x\le837.27+\lambda/2+1.9608\,s+2.9715\,s\lambda$. From $\log(s/1000)\le s/1000-1$ and $\log1000=3\log2+3\log5\le6.9078$ we get $\lambda=2\log s\le11.8157+s/500$. Substituting this and $x=s^2$ yields $0.0478\,s^2-37.08\,s-843.2\le0$, which is false because $s^2\ge1000s$ (one \texttt{linarith} call in Lean).
\end{proof}

\subsection{The Bertrand chain}

\begin{lemma}[\decl{chain_step}, \decl{bertrand_ge}]\label{lem:chain}
For $k\ge8$: $B_k\ge83$, $5B_k<3B_{k+1}$, $B_{k+1}<2B_k$, and no prime lies in $(B_{k+1},2B_k)$.
\end{lemma}

\begin{proof}
If $b\ge83$ and $c$ is the largest prime $\le2b-1$, Proposition~\ref{prop:nagura} gives a prime $p\le2b-1$ with $5b<3p$, so $5b<3c$; the other claims hold by the choice of $c$. Now induct from $B_8=83$ (a kernel computation).
\end{proof}

\subsection{The factorial divisibility}

\begin{lemma}[\decl{pointwise}]\label{lem:pointwise}
Let $a<2b$ and $5a<3c$. For an integer $y\ge1$ put $k=\fl{2a/y}$ and $f(y)=\fl{2b/y}+\fl{2c/y}-2k$. Then \textup{(i)} $f(y)\ge-1$; \textup{(ii)} $f(y)\ge0$ unless $c<y\le2a$; \textup{(iii)} $f(y)\ge1$ if $k\ge13$.
\end{lemma}

\begin{proof}
Let $q_1=\fl{2b/y}$, $q_2=\fl{2c/y}$. From $ky\le2a<4b<2(q_1+1)y$ and $5ky\le10a<6c<3(q_2+1)y$ we get $2q_1\ge k-1$ and $3q_2\ge5k-2$, so $q_1+q_2\ge(13k-7)/6$, which is $\ge2k-1$ for $k\ge1$ and $\ge2k+1$ for $k\ge13$. This gives (iii), and (i) since $f\ge0$ when $k=0$. For (ii): $y>2a$ forces $k=0$. If $y\le c$, then $q_2\ge2$, which gives $f\ge0$ for $k\le1$; for $k\ge2$ the bounds $q_1\ge\lceil(k-1)/2\rceil$ and $q_2\ge\lceil(5k-2)/3\rceil$ give $q_1+q_2\ge2k$ (directly for $2\le k\le6$, and from $(13k-7)/6\ge2k$ for $k\ge7$). In Lean all cases are closed by \texttt{omega}.
\end{proof}

\begin{lemma}[\decl{fact_dvd}]\label{lem:factdvd}
Let $b\ge83$, $5b<3a$, $a<2b$, $5a<3c$, $c<2a$, and suppose that no prime lies in $(c,2a)$. Then $(2a)!^2\mid(2b)!\,(2c)!$.
\end{lemma}

\begin{proof}
By Legendre's formula (Mathlib's \decl{Nat.factorization_factorial}) it suffices to show $\sum_{j\ge1}f(\ell^j)\ge0$ for every prime $\ell$, with $f$ as in Lemma~\ref{lem:pointwise}. If no power $\ell^j$ ($j\ge1$) lies in $(c,2a]$, every term is $\ge0$ by Lemma~\ref{lem:pointwise}(ii). Otherwise let $c<\ell^j\le2a$. This $j$ is unique: if also $c<\ell^i\le2a$ with $i<j$, then $2a\ge\ell^j\ge2\ell^i>2c$, contradicting $c>5a/3$. Moreover $j\ge2$: a prime $\ell\in(c,2a)$ is excluded by hypothesis, and $\ell=2a$ is impossible since $2a>2$ is even. Hence $\ell^2\le2a$, and so $13\ell\le2a$ (from $\ell^2\le2a$ if $\ell\ge13$, and from $2a>10b/3>156$ if $\ell<13$). Thus $\fl{2a/\ell}\ge13$, so $f(\ell)\ge1$ by (iii), $f(\ell^j)\ge-1$ by (i), and all other terms are $\ge0$ by (ii).
\end{proof}

\subsection{The rows}

\begin{lemma}[\decl{step_core}]\label{lem:step}
Under the hypotheses of Lemma~\ref{lem:factdvd}, if $c$ is prime, then $2c-2a$ is the least $k\ge2a-2b$ with $\rf{(2b+1)}{2a-2b}\mid\rf{(2a+1)}{k}$.
\end{lemma}

\begin{proof}
First, $c>5a/3$ and $b>a/2$ give $b+c>2a$, i.e.\ $2c-2a\ge2a-2b$. Next, $(2b)!\,\rf{(2b+1)}{2a-2b}=(2a)!$ and $(2a)!\,\rf{(2a+1)}{2c-2a}=(2c)!$; multiplying the desired divisibility by $(2b)!\,(2a)!$ turns it into Lemma~\ref{lem:factdvd}, so $k=2c-2a$ is admissible. Finally, $c>5a/3>25b/9>2b$ and $c<2a$, so the prime $c$ divides $\rf{(2b+1)}{2a-2b}$. If $k<2c-2a$, the factors $2a+1,\dots,2a+k$ of $\rf{(2a+1)}{k}$ lie strictly between $c$ and $2c$, so $c\nmid\rf{(2a+1)}{k}$ and $k$ is not admissible. (This part does not use the bound $k\ge2a-2b$.)
\end{proof}

\begin{proposition}[\decl{rows}]\label{prop:rows}
For every $n\ge10$, $a(n)=2B_{n-2}+1$ and $L(n)=2B_{n-1}-2B_{n-2}$.
\end{proposition}

\begin{proof}
Induction on $n$. For $n=10$ the kernel evaluates $a(10)=167$, $L(10)=160$, $B_8=83$ and $B_9=163$ (\texttt{decide +kernel}). For the step put $b=B_{n-2}$, $a=B_{n-1}$, $c=B_n$. Lemma~\ref{lem:chain} gives the hypotheses of Lemmas~\ref{lem:factdvd} and~\ref{lem:step}. By the induction hypothesis row $n$ is $[2b+1,2a]$, so $a(n+1)=2a+1$, and \decl{len_succ_isLeast} together with Lemma~\ref{lem:step} gives $L(n+1)=2c-2a$.
\end{proof}

\subsection{Proof of Theorem~\ref{thm:main}}
For $n=5,\dots,10$, row $n-1$ is $[5,8]$, $[9,14]$, $[15,26]$, $[27,46]$, $[47,86]$, $[87,166]$, with largest prime $p=7,13,23,43,83,163$, which is $B_{n-1}$, and $e(n)=14,26,46,86,166,326=2p$. These values are evaluated by the kernel (\texttt{decide +kernel}), and the absence of primes in $(p,e(n-1)]$ is checked by \texttt{interval\_cases} and \texttt{norm\_num} (\decl{small_case}). Let $n\ge11$. By Proposition~\ref{prop:rows}, row $n-1$ is $[2B_{n-3}+1,2B_{n-2}]$ and $e(n)=2B_{n-1}$. By the chain inequalities, $p=B_{n-1}$ satisfies $p>25B_{n-3}/9>2B_{n-3}$ and $p<2B_{n-2}$, so it is a prime in row $n-1$. A prime $q>p$ in row $n-1$ would satisfy $q<2B_{n-2}$ (as $2B_{n-2}>2$ is even), contradicting the absence of primes in $(B_{n-1},2B_{n-2})$. Hence $p$ is the largest prime of row $n-1$, and $e(n)=2p$. \qed

\section{Machine verification}\label{sec:lean}

\subsection*{Files and audit}
The proof is the file \path{lean/Math0/Frontier/OEIS/A090905.lean} of the math0 repository, \url{https://github.com/Foundation-0/math0}. The main declaration is \decl{Math0.Frontier.A090905.result}, of type \decl{Math0.Statements.A090905.statement}; all other declarations in the file are private. The file uses no \texttt{sorry}, \texttt{native\_decide}, \texttt{axiom} or \texttt{set\_option}; the finite computations use \texttt{decide}, \texttt{decide +kernel} and \texttt{norm\_num}, whose results are checked by the kernel. The file ends with \lstinline|#math0_settles result statement| (with fully qualified names), a command defined in \path{lean/Math0/Audit.lean} for statements formalized by math0. It fails unless \decl{statement} is a \texttt{Prop}-valued definition without universe parameters, the type of \decl{result} is syntactically \decl{statement} or its negation (up to metadata, without unfolding), and both depend only on the axioms \decl{propext}, \decl{Classical.choice} and \decl{Quot.sound} (which excludes \decl{sorryAx}, \texttt{native\_decide} and custom axioms). It reports:
\begin{lstlisting}[basicstyle=\ttfamily\scriptsize,belowskip=0pt]
math0 audit: 'Math0.Statements.A090905.statement' depends on axioms [Classical.choice, Quot.sound, propext]
math0 audit: 'Math0.Frontier.A090905.result' depends on axioms [Classical.choice, Quot.sound, propext]
math0 statement check: 'Math0.Frontier.A090905.result' proves 'Math0.Statements.A090905.statement'
\end{lstlisting}

\subsection*{Toolchain and reproduction}
We use Lean~4 \cite{Lean4} \texttt{v4.33.1} and Mathlib \cite{Mathlib} \texttt{v4.33.1}, pinned in \path{lean/}; the repository also pins formal-conjectures \cite{FC} at commit \texttt{c77fdd4}, which this result does not use. To reproduce, run \path{scripts/verify.sh}: it fetches the Mathlib cache, builds the library with \texttt{lake build} (re-running every audit), checks that the pinned \texttt{formal\_conjectures} and \texttt{mathlib} checkouts are unmodified, and counts the passed statement checks.

{\footnotesize\raggedright
\begin{thebibliography}{99}

\bibitem{Cheb}
P.~L. Tchebichef, \emph{M\'emoire sur les nombres premiers}, J. Math.\ Pures Appl.\ (1) \textbf{17} (1852), 366--390.

\bibitem{FC}
M.~Firsching, P.~Lezeau, S.~Mercuri, M.~Z. Horv\'ath, Y.~Dillies, C.~S\"onne, E.~Wieser, F.~Zhang, T.~Hubert, B.~Ag\"uera y Arcas, and P.~Kohli, \emph{Formal Conjectures: An open and evolving benchmark for verified discovery in mathematics}, arXiv:2605.13171, 2026. Repository: \url{https://github.com/google-deepmind/formal-conjectures}.

\bibitem{Mathlib}
The mathlib Community, \emph{The Lean mathematical library}, in: Proceedings of the 9th ACM SIGPLAN International Conference on Certified Programs and Proofs (CPP 2020), ACM, 2020, 367--381.

\bibitem{Lean4}
L.~de Moura and S.~Ullrich, \emph{The Lean 4 theorem prover and programming language}, in: Automated Deduction -- CADE 28, Lecture Notes in Comput.\ Sci.\ \textbf{12699}, Springer, 2021, 625--635.

\bibitem{Nag}
J.~Nagura, \emph{On the interval containing at least one prime number}, Proc.\ Japan Acad.\ \textbf{28} (1952), 177--181.

\bibitem{A006992}
The OEIS Foundation Inc., A006992 (Bertrand primes), \url{https://oeis.org/A006992}.

\bibitem{A090904}
The OEIS Foundation Inc., A090904, A090906, A090907 (A.~Murthy, 2003), \url{https://oeis.org/A090904}, \url{https://oeis.org/A090906}, \url{https://oeis.org/A090907}.

\bibitem{A090905}
The OEIS Foundation Inc., A090905 (A.~Murthy, Dec.~13, 2003), \url{https://oeis.org/A090905}.

\bibitem{A093911}
The OEIS Foundation Inc., A093911 (A.~Murthy, 2004), A113117, A113118 (L.~Quet, 2006), \url{https://oeis.org/A093911}, \url{https://oeis.org/A113117}, \url{https://oeis.org/A113118}.

\bibitem{A381901}
The OEIS Foundation Inc., A381901 (A.~Niedermaier, Mar.~9, 2025), \url{https://oeis.org/A381901}.

\end{thebibliography}
}

\end{document}
